\documentclass[10pt]{extarticle}
\usepackage[a4paper,margin=1.5cm]{geometry}
\usepackage[T1]{fontenc}
\usepackage{amsmath,amssymb}
\usepackage{graphicx}
\usepackage[export]{adjustbox}
\usepackage{multicol}
\usepackage{enumitem}
\usepackage{booktabs,array,tabularx}
\setlength{\columnsep}{0.7cm}
\setlength{\parindent}{0pt}
\setlength{\parskip}{3pt}
\pagestyle{empty}
\makeatletter
% figure registry (figdefs.tex is generated by the extract step)
\newcommand{\qfigdef}[4]{\@namedef{qf@file@#1}{#2}\@namedef{qf@w@#1}{#3}\@namedef{qf@wide@#1}{#4}}
\newcommand{\qf@check}[1]{\@ifundefined{qf@file@#1}{\PackageError{score2}{Unknown figure id #1}{}}{}}
\newcommand{\qf@img}[4][]{\includegraphics[width=#2\dimexpr\@nameuse{qf@w@#4}pt\relax,max width=#3,max height=0.45\textheight,#1]{\@nameuse{qf@file@#4}}}
% Full-width questions: finalize wraps every question that contains a wide figure in \qfullbegin...\qfullend,
% so the WHOLE question (stem, figures, options) is set across the page in reading order, then the two columns
% resume. The question's head - its number, stem and first figure (\qfullsplit marks the end of the head) -
% is kept on one page, so a number is never left at a page foot with its figure overleaf; the rest may break
% across pages normally, so no big gaps are left. Without a \qfullsplit, the whole question is kept together
% if it is under 45% of a page. Never write these yourself.
\newif\ifqf@full
\newif\ifqf@boxed
\newsavebox{\qf@box}
\newcommand{\qf@open}{\global\qf@fulltrue\par\global\qf@boxedtrue
  \setbox\qf@box\vbox\bgroup\hsize\textwidth\linewidth\textwidth}
\newcommand{\qf@emit}[1]{\par\egroup\global\qf@boxedfalse
  \ifdim\dimexpr\ht\qf@box+\dp\qf@box\relax<#1\textheight\noindent\usebox{\qf@box}\par\else\unvbox\qf@box\fi}
\newcommand{\qfullsplit}{\ifqf@boxed\qf@emit{0.85}\fi}
\newcommand{\qf@close}{\ifqf@boxed\qf@emit{0.45}\fi\global\qf@fullfalse}
\newcommand{\qfullbegin}{\end{multicols}\qf@open}
\newcommand{\qfullend}{\qf@close\begin{multicols}{2}}
\newcommand{\qfullstart}{\qf@open}
\newcommand{\qfulldone}{\qf@close}
% \qfig[scale]{id}: block figure, centred, directly where it is written (size set by the tool).
% Fallback if a wide figure is not inside a full-width question: break out of the columns in place - no floats.
\newcommand{\qfig}[2][1]{\qf@check{#2}%
  \ifqf@full\qf@blk{#1}{#2}\else
  \ifnum\@nameuse{qf@wide@#2}=1 \expandafter\@firstoftwo\else\expandafter\@secondoftwo\fi
  {\end{multicols}\begin{center}\qf@img{#1}{\textwidth}{#2}\end{center}\begin{multicols}{2}}%
  {\qf@blk{#1}{#2}}\fi}
\newcommand{\qf@blk}[2]{\par\smallskip{\centering\qf@img{#1}{\linewidth}{#2}\par}\smallskip}
% \qfigi[scale]{id}: inline figure (inside option lists / tables / minipages); never breaks columns.
\newcommand{\qfigi}[2][1]{\qf@check{#2}\qf@img[valign=c]{#1}{\linewidth}{#2}}
\makeatother
% structure macros - lint reads these to check completeness and figure placement
\newcommand{\papersec}[1]{\par\bigskip{\centering\Large\bfseries PAPER-#1\par}\medskip}
\newcommand{\partsec}[2]{\par\medskip\noindent{\large\bfseries PART-#1: #2}\par\smallskip}
\newcommand{\sect}[1]{\par\smallskip\noindent\textbf{#1}\par}
\newcommand{\qn}[1]{\par\smallskip\noindent\textbf{#1.}~}
\newcommand{\paraq}[2]{\par\smallskip\noindent\textbf{Paragraph for Questions #1 and #2}\par}
\newenvironment{qoptions}{\begin{itemize}[leftmargin=1.8em,itemsep=1pt,topsep=2pt,label={}]}{\end{itemize}}
\newcommand{\opt}[2]{\item[(#1)] #2}
% figure registry (generated by score2_tools.py finalize)
\qfigdef{p1_phy_s1_q05_1}{TOAS - Test 19 (05 Feb 2023)-img/p1_phy_s1_q05_1.png}{74.0}{0}
\qfigdef{p1_phy_s1_q06_1}{TOAS - Test 19 (05 Feb 2023)-img/p1_phy_s1_q06_1.png}{137.6}{0}
\qfigdef{p1_phy_s1_q08_1}{TOAS - Test 19 (05 Feb 2023)-img/p1_phy_s1_q08_1.png}{89.9}{0}
\qfigdef{p1_phy_s2_q09_1}{TOAS - Test 19 (05 Feb 2023)-img/p1_phy_s2_q09_1.png}{150.8}{0}
\qfigdef{p1_phy_s2_q10_1}{TOAS - Test 19 (05 Feb 2023)-img/p1_phy_s2_q10_1.png}{150.8}{0}
\qfigdef{p1_phy_s2_q11_1}{TOAS - Test 19 (05 Feb 2023)-img/p1_phy_s2_q11_1.png}{108.3}{0}
\qfigdef{p1_phy_s2_q12_1}{TOAS - Test 19 (05 Feb 2023)-img/p1_phy_s2_q12_1.png}{134.7}{0}
\qfigdef{p1_phy_s3_q01_1}{TOAS - Test 19 (05 Feb 2023)-img/p1_phy_s3_q01_1.png}{82.0}{0}
\qfigdef{p1_phy_s3_q02_1}{TOAS - Test 19 (05 Feb 2023)-img/p1_phy_s3_q02_1.png}{79.4}{0}
\qfigdef{p1_phy_s3_q03_1}{TOAS - Test 19 (05 Feb 2023)-img/p1_phy_s3_q03_1.png}{103.2}{0}
\qfigdef{p1_phy_s3_q04_1}{TOAS - Test 19 (05 Feb 2023)-img/p1_phy_s3_q04_1.png}{58.2}{0}
\qfigdef{p1_phy_s3_q05_1}{TOAS - Test 19 (05 Feb 2023)-img/p1_phy_s3_q05_1.png}{108.2}{0}
\qfigdef{p1_chem_s1_q05_A}{TOAS - Test 19 (05 Feb 2023)-img/p1_chem_s1_q05_A.png}{89.7}{0}
\qfigdef{p1_chem_s1_q05_B}{TOAS - Test 19 (05 Feb 2023)-img/p1_chem_s1_q05_B.png}{112.3}{0}
\qfigdef{p1_chem_s1_q05_C}{TOAS - Test 19 (05 Feb 2023)-img/p1_chem_s1_q05_C.png}{79.4}{0}
\qfigdef{p1_chem_s1_q05_D}{TOAS - Test 19 (05 Feb 2023)-img/p1_chem_s1_q05_D.png}{71.3}{0}
\qfigdef{p1_chem_s1_q06_A}{TOAS - Test 19 (05 Feb 2023)-img/p1_chem_s1_q06_A.png}{65.6}{0}
\qfigdef{p1_chem_s1_q06_B}{TOAS - Test 19 (05 Feb 2023)-img/p1_chem_s1_q06_B.png}{70.3}{0}
\qfigdef{p1_chem_s1_q06_C}{TOAS - Test 19 (05 Feb 2023)-img/p1_chem_s1_q06_C.png}{70.8}{0}
\qfigdef{p1_chem_s1_q06_D}{TOAS - Test 19 (05 Feb 2023)-img/p1_chem_s1_q06_D.png}{70.3}{0}
\qfigdef{p1_chem_s1_q07_A}{TOAS - Test 19 (05 Feb 2023)-img/p1_chem_s1_q07_A.png}{95.1}{0}
\qfigdef{p1_chem_s1_q07_B}{TOAS - Test 19 (05 Feb 2023)-img/p1_chem_s1_q07_B.png}{87.4}{0}
\qfigdef{p1_chem_s1_q07_C}{TOAS - Test 19 (05 Feb 2023)-img/p1_chem_s1_q07_C.png}{89.7}{0}
\qfigdef{p1_chem_s1_q07_D}{TOAS - Test 19 (05 Feb 2023)-img/p1_chem_s1_q07_D.png}{92.5}{0}
\qfigdef{p1_chem_s1_q08_B}{TOAS - Test 19 (05 Feb 2023)-img/p1_chem_s1_q08_B.png}{118.3}{0}
\qfigdef{p1_chem_s1_q08_C}{TOAS - Test 19 (05 Feb 2023)-img/p1_chem_s1_q08_C.png}{213.3}{0}
\qfigdef{p1_chem_s1_q08_D}{TOAS - Test 19 (05 Feb 2023)-img/p1_chem_s1_q08_D.png}{189.6}{0}
\raggedcolumns
\makeatletter\@beginparpenalty=10000\makeatother
\newcolumntype{L}[1]{>{\raggedright\arraybackslash}p{#1}}
\begin{document}
\begin{center}{\Large\bfseries TOAS --- Test 19 (05 Feb 2023)}\end{center}
\begin{multicols}{2}
\raggedcolumns
\papersec{1}
\partsec{1}{PHYSICS}
\sect{SECTION-I(i)}
\qn{1} On a smooth horizontal surface there is a disc of mass $2m$ and radius $R$ is placed. An insect of mass $m$ sits gently at a distance $\dfrac R2$ from center of disc and crawls uniformly with speed `$v_0$' with respect to centre of disc along circular path with respect to centre of disc.
\begin{qoptions}
\opt{A}{With respect to a ground observer insect moves in circular path of radius $\dfrac R2$.}
\opt{B}{With respect to a ground observer centre of disc moves on a circular path of radius $\dfrac R6$.}
\opt{C}{Net external force on this system is zero.}
\opt{D}{Insect applies a force $\dfrac4{27}mv_0^2R$ on disc tangentially to the surface of disc.}
\end{qoptions}
\qn{2} Choose the \textit{INCORRECT} statement(s) :-
\begin{qoptions}
\opt{A}{All the system in stable equilibrium, if disturbed perform simple harmonic motion.}
\opt{B}{All the system in unstable equilibrium, if disturbed do not perform simple harmonic motion.}
\opt{C}{All the system in stable equilibrium, if disturbed perform oscillatory motion.}
\opt{D}{Any system in unstable equilibrium, if disturbed will not perform oscillatory motion.}
\end{qoptions}
\qn{3} An ideal monoatomic gas is taken through a thermodynamics process $P=aV+b$ when $P$ and $V$ pressure and volume of gas and $a$ and $b$ are positive constants.
\begin{qoptions}
\opt{A}{If $b=0$, molar heat capacity of gas is constant.}
\opt{B}{If $a\neq0$, $b\neq0$, molar heat capacity of gas is lower in higher volume}
\opt{C}{If $a=0$, increase in volume per unit rise in temperature is constant.}
\opt{D}{Slope of temperature versus volume graph is constant.}
\end{qoptions}
\qn{4} A particle starting from rest moves with time varying acceleration $\vec a=3\hat i+4t\hat j$ m/s$^2$
\begin{qoptions}
\opt{A}{Speed of particle after 1 sec should be 5 m/s.}
\opt{B}{At $t=\dfrac32$ sec particle will be moving at $45^\circ$ with $x$ axis.}
\opt{C}{At $t=1$ sec rate of change of speed will be $\dfrac{17}{\sqrt{13}}$ m/s$^2$}
\opt{D}{At $t=1$ sec radius of curvature of the path will be $\dfrac{13}6\sqrt{13}$ m}
\end{qoptions}
\qn{5} Figure shows three leg table, with top as a uniform equilateral triangular board of mass 10 kg. Two bodies of masses 2 kg and 3 kg are placed at corners A and C. The table is placed on a horizontal surface.
\qfig{p1_phy_s1_q05_1}
\begin{qoptions}
\opt{A}{Normal force on leg at A $=\dfrac{160}3$ N}
\opt{B}{Normal force on leg at B $=\dfrac{100}3$ N}
\opt{C}{Normal force on leg at C $=\dfrac{190}3$ N}
\opt{D}{If both the blocks are removed, normal force on leg at B will remain same.}
\end{qoptions}
\qn{6} ABC is a uniform chain of mass 10 kg with B as lowest point.
\qfig{p1_phy_s1_q06_1}
\begin{qoptions}
\opt{A}{Tension in chain C $>$ tension at A}
\opt{B}{Tension at B $=48$ N}
\opt{C}{Mass ratio of AB and BC $=9:16$}
\opt{D}{Mass of protion AB is 3.6 kg}
\end{qoptions}
\qn{7} 20 gm ice is kept in a calorimeter at $0^\circ$C. Now 40 gm water at $30^\circ$C is poured gently into it. There is no leakage of heat. In equilibrium.
\begin{qoptions}
\opt{A}{Final temperature will be $0^\circ$C only if heat capacity of calorimeter is neglected.}
\opt{B}{Final temperature will be $0^\circ$C irrespective of value of heat capacity of calorimeter.}
\opt{C}{75\% of ice will melt.}
\opt{D}{Remaining ice will be more than 10 gm.}
\end{qoptions}
\qn{8} Figure shows a smooth circular groove in vertical plane in which a uniform rod of length $\ell=1.2\,R$ is released from the position shown.
\qfig{p1_phy_s1_q08_1}
\begin{qoptions}
\opt{A}{Initial acceleration of centre of mass of rod is $\dfrac{16}{19}g$.}
\opt{B}{Maximum speed of centre of mass of rod will be $\sqrt{1.6gR}$}
\opt{C}{Maximum angular speed of rod will be $\sqrt{\dfrac{40}{19}\dfrac gR}$}
\opt{D}{Normal force at one end of rod when it becomes horizontal is $mg$.}
\end{qoptions}
\sect{SECTION-I(ii)}
\qn{9} Figure shows a rider performing a stunt with rear wheel in pure rolling.

$M=$ total mass of system $=200$ kg,

$m=$ mass of rear wheel (ring) $=10$ kg,

$R=$ radius of wheel $=50$ cm,

$h=$ height of centre of mass of system,

$b=$ distance of vertical line through centre of mass from contact point of rear wheel.
\qfig{p1_phy_s2_q09_1}
\begin{center}\footnotesize\setlength{\tabcolsep}{2pt}
\begin{tabular}{|L{0.8cm}L{3.6cm}|L{0.6cm}L{1.0cm}|}\hline
\multicolumn{2}{|c|}{List-I}&\multicolumn{2}{c|}{List-II}\\\hline
(I)&Frictional force in KN&(P)&1.2\\\hline
(II)&Torque provided by engine in KN-m&(Q)&0.630\\\hline
(III)&Value of $h$ in meter&(R)&50\\\hline
(IV)&Value of $b$ in centimeter&(S)&1\\\hline
&&(T)&60\\\hline
&&(U)&0.525\\\hline
\end{tabular}\end{center}
It the motorcycle in the given situation accelerates with 5 m/s$^2$, the correct matching is as :-
\begin{qoptions}
\opt{A}{I $\to$ S;II $\to$ U;III $\to$ S;IV $\to$ R}
\opt{B}{I $\to$ S;II $\to$ U;III $\to$ R;IV $\to$ S}
\opt{C}{I $\to$ S;II $\to$ Q;III $\to$ S;IV $\to$ S}
\opt{D}{I $\to$ P;II $\to$ Q;III $\to$ R;IV $\to$ S}
\end{qoptions}
\qn{10} Figure shows a rider performing a stunt with rear wheel in pure rolling.

$M=$ total man of system $=200$ kg,

$m=$ mass of rear wheel (ring) $=10$ kg,

$R=$ radius of wheel $=50$ cm,

$h=$ height of centre of mass of system,

$b=$ distance of vertical line through centre of mass from contact point of rear wheel.
\qfig{p1_phy_s2_q10_1}
\begin{center}\footnotesize\setlength{\tabcolsep}{2pt}
\begin{tabular}{|L{0.8cm}L{3.6cm}|L{0.6cm}L{1.0cm}|}\hline
\multicolumn{2}{|c|}{List-I}&\multicolumn{2}{c|}{List-II}\\\hline
(I)&Frictional force in KN&(P)&1.2\\\hline
(II)&Torque provided by engine in KN-m&(Q)&0.630\\\hline
(III)&Value of $h$ in meter&(R)&50\\\hline
(IV)&Value of $b$ in centimeter&(S)&1\\\hline
&&(T)&60\\\hline
&&(U)&0.525\\\hline
\end{tabular}\end{center}
It the motorcycle in the given situation accelerates with 6 m/s$^2$, the correct matching is as :-
\begin{qoptions}
\opt{A}{I $\to$ P;II $\to$ Q;III $\to$ S;IV $\to$ T}
\opt{B}{I $\to$ P;II $\to$ Q;III $\to$ T;IV $\to$ S}
\opt{C}{I $\to$ P;II $\to$ S;III $\to$ R;IV $\to$ Q}
\opt{D}{I $\to$ R;II $\to$ S;III $\to$ T;IV $\to$ U}
\end{qoptions}
\qn{11} A mass of 1kg is supported by a cord 5m long attached to a massless ring free to move on a horizontal fixed rod, as shown in the figure. Coefficient of friction between ring and the rod is 0.75. A second cord is attached to the weight and passes through a pulley at $\dfrac{25}3$m from ring. Block A of mass $m$ is attached to the other end of this cord. Just after released.
\qfig{p1_phy_s2_q11_1}
\begin{center}\footnotesize\setlength{\tabcolsep}{2pt}
\begin{tabular}{|L{0.8cm}L{3.9cm}|L{0.6cm}L{0.8cm}|}\hline
\multicolumn{2}{|c|}{List-I}&\multicolumn{2}{c|}{List-II}\\\hline
(I)&Tension in the string connected to block A in `N'.&(P)&6\\\hline
(II)&Tension in the 5m cord in `N'&(Q)&8\\\hline
(III)&The value of $\theta$ is $\dfrac\pi k$ radian, possible value of $k$.&(R)&2\\\hline
(IV)&Acceleration of block A in m/s$^2$&(S)&0\\\hline
&&(T)&4\\\hline
&&(U)&12\\\hline
\end{tabular}\end{center}
If the ring is just about to slip
\begin{qoptions}
\opt{A}{I $\to$ P;II $\to$ Q;III $\to$ R;IV $\to$ S}
\opt{B}{I $\to$ Q;II $\to$ P;III $\to$ T;IV $\to$ S}
\opt{C}{I $\to$ T;II $\to$ Q;III $\to$ R;IV $\to$ S}
\opt{D}{I $\to$ U;II $\to$ Q;III $\to$ R;IV $\to$ S}
\end{qoptions}
\qn{12} A mass of 1kg is supported by a cord 5m long attached to a massless ring free to move on a horizontal fixed rod, as shown in the figure. Coefficient of friction between ring and the rod is 0.75. A second cord is attached to the weight and passes through a pulley at $\dfrac{25}3$m from ring. Block A of mass $m$ is attached to the other end of this cord. Just after released.
\qfig{p1_phy_s2_q12_1}
\begin{center}\footnotesize\setlength{\tabcolsep}{2pt}
\begin{tabular}{|L{0.8cm}L{3.9cm}|L{0.6cm}L{0.8cm}|}\hline
\multicolumn{2}{|c|}{List-I}&\multicolumn{2}{c|}{List-II}\\\hline
(I)&Tension in the string connected to block A in `N'.&(P)&6\\\hline
(II)&Tension in the 5m cord in `N'&(Q)&8\\\hline
(III)&The value of $\theta$ is $\dfrac\pi k$ radian, possible value of $k$.&(R)&2\\\hline
(IV)&Acceleration of block A in m/s$^2$&(S)&0\\\hline
&&(T)&4\\\hline
&&(U)&12\\\hline
\end{tabular}\end{center}
The value of $m$ was initially chosen such that the ring just began to slide. Now the ring is fixed to the rod and mass of block A is increased to 5m, holding the 1kg mass. Just after the system is released, then
\begin{qoptions}
\opt{A}{I $\to$ P;II $\to$ Q;III $\to$ R;IV $\to$ T}
\opt{B}{I $\to$ Q;II $\to$ P;III $\to$ T;IV $\to$ U}
\opt{C}{I $\to$ T;II $\to$ U;III $\to$ R;IV $\to$ Q}
\opt{D}{I $\to$ U;II $\to$ Q;III $\to$ R;IV $\to$ P}
\end{qoptions}
\sect{SECTION-II}
\qn{1} A mass 1kg moves initially with a velocity $v_0=10\sqrt{15}$ m/s. A constant power $P=1000$ W is applied so as to increase velocity find the time in seconds after which acceleration is reduced to half of its original value.
\qfig{p1_phy_s3_q01_1}
\qn{2} A small spherical ball of density equal to density of liquid is released from rest in a vessel filled completely with a liquid and accelerating with acceleration $5\hat i+5\hat j$ m/s$^2$. Find initial acceleration of the ball just after the release with respect to vessel.
\qfig{p1_phy_s3_q02_1}
\qn{3} An old record player of 10 cm radius turns at 10 rad/sec while mounted on a $30^\circ$ incline as in the figure. A particle of mass $m$ can be placed any where on the rotating record. If the least possible coefficient of friction $\mu$ that must exist for no slipping to occur is `$\mu$'. Find $2\sqrt3\mu$.
\qfig{p1_phy_s3_q03_1}
\qn{4} The ends of a chain lie in piles at A and C of negligible height. When given an initial speed 10 m/s, the chain keeps on moving at that speed over the fixed pulley B. Neglect friction, Find $h$ in meter. Assume chain do not tangle.
\qfig{p1_phy_s3_q04_1}
\qn{5} A black surface at constant high temperature $T_H$ is parallel to another black plane surface at a constant lower temperature $T_L$ (enclosed volume in between). A heat shield consisting of two thin black plates is placed between the warm and the cold surfaces and parallel to these. After some time steady state is achieved. The ratio of initial heat current to final heat current is (Assume all surfaces to be large with respect to distance between surfaces)
\qfig{p1_phy_s3_q05_1}
\qn{6} A man crosses a river of width $d$. Flow speed of water is $v$. Speed of swimmer relative to water is also $v$. Man always heads towards the point exactly opposite to the starting point at the other bank (relative to water). If total drift is $\dfrac{md}n$. Find minimum possible value of $m+n$. ($m$ \& $n$ are integers.)
\partsec{2}{CHEMISTRY}
\sect{SECTION-I(i)}
\qn{1} In a closed rigid 2 L container the following reaction is occuring :

$\mathrm{PCl_5(g)\rightleftharpoons PCl_3(g)+Cl_2(g)}$

If initially taken 2 moles of PCl$_5$ dissociates to 60\% to establish equilibrium then which of the following is/are correct statement(s) :
\begin{qoptions}
\opt{A}{Mass inside the container remains same.}
\opt{B}{Moles inside the container increases.}
\opt{C}{Molar mass of equilibrium mixture is lesser than molar mass of PCl$_5$.}
\opt{D}{$K_C<1$}
\end{qoptions}
\qn{2} One of the major use of H$_2$O$_2$ is for artificial respiration. It is stored in cylinder of 3 litre capacity completely filled with H$_2$O$_2$ solution of `45.4V'. In a hospital of kota, patient under artificial respiration takes 200 ml O$_2$ per min at 1 atm at 273 K for 2.8 hours from cylinder.

[Assume volume of solution and rate of decomposition remain constant]

Then which of the following statement(s) is/are correct.
\begin{qoptions}
\opt{A}{Initial moles of H$_2$O$_2$ in the cylinder are 12.}
\opt{B}{Volume of oxygen used for respiration is 33.6 L.}
\opt{C}{Volume strength of H$_2$O$_2$ left in solution is 11.35 V.}
\opt{D}{Volume strength of H$_2$O$_2$ left in solution is 34.05 V.}
\end{qoptions}
\qn{3} Select the correct order(s) with respect to specified property.
\begin{qoptions}
\opt{A}{NaCl $>$ KCl (solubility in water)}
\opt{B}{LiH $>$ NaH (Thermal stability)}
\opt{C}{N(CH$_3$)$_3$ $>$ N(SiH$_3$)$_3$ (lewis basic strength)}
\opt{D}{(CH$_3$)$_3$SiOH $<$ (CH$_3$)$_3$COH (H$^+$ donating ability)}
\end{qoptions}
\qn{4} Select the incorrect statement(s) regarding following hydrolysis processes :-
\begin{qoptions}
\opt{A}{In hydrolysis of NCl$_3$ , as nitrogen is central atom, HNO$_3$ and HCl is formed}
\opt{B}{Hydrolysis of XeF$_2$ and XeF$_4$ is redox change whereas hydrolysis of XeF$_6$ is not}
\opt{C}{During the stepwise hydrolysis of P$_4$O$_{10}$, various oxyacids of phosphorous such as meta, poly, pyro and hypo acids are formed}
\opt{D}{The number of BF$_3$ molecules hydrolysed per mole H$_2$O is higher than the number of SiF$_4$ molecules hydrolysed per mole H$_2$O}
\end{qoptions}
\qn{5} Identify the structure of transdecaline
\begin{qoptions}
\opt{A}{\qfigi{p1_chem_s1_q05_A}}
\opt{B}{\qfigi{p1_chem_s1_q05_B}}
\opt{C}{\qfigi{p1_chem_s1_q05_C}}
\opt{D}{\qfigi{p1_chem_s1_q05_D}}
\end{qoptions}
\qn{6} Identify the structure of identical compounds
\begin{qoptions}
\opt{A}{\qfigi{p1_chem_s1_q06_A}}
\opt{B}{\qfigi{p1_chem_s1_q06_B}}
\opt{C}{\qfigi{p1_chem_s1_q06_C}}
\opt{D}{\qfigi{p1_chem_s1_q06_D}}
\end{qoptions}
\qn{7} Identify the geometrical isomers
\begin{qoptions}
\opt{A}{\qfigi{p1_chem_s1_q07_A}}
\opt{B}{\qfigi{p1_chem_s1_q07_B}}
\opt{C}{\qfigi{p1_chem_s1_q07_C}}
\opt{D}{\qfigi{p1_chem_s1_q07_D}}
\end{qoptions}
\qn{8} Identify correct statement?
\begin{qoptions}
\opt{A}{F$^\ominus<$ Cl$^\ominus<$ Br$^\ominus<$ I$^\ominus$ Nucleophilicity order in H$_2$O}
\opt{B}{\qfigi{p1_chem_s1_q08_B}}
\opt{C}{\qfigi{p1_chem_s1_q08_C}}
\opt{D}{\qfigi{p1_chem_s1_q08_D}}
\end{qoptions}
\sect{SECTION-I(ii)}
\qn{9} \begin{center}\footnotesize\setlength{\tabcolsep}{2pt}
\begin{tabular}{|L{0.8cm}L{2.2cm}|L{0.6cm}L{2.2cm}|}\hline
\multicolumn{2}{|L{3.0cm}|}{\centering List-I\\(condition of vander walls gas)}&\multicolumn{2}{L{2.8cm}|}{\centering List-II\\(compressibility factor)}\\\hline
(I)&At low to moderate pressure&(P)&$Z=1+\dfrac{Pb}{RT}$\\\hline
(II)&At critical condition of $P,V,T$&(Q)&$Z=1-\dfrac{a}{V_mRT}$\\\hline
(III)&For H$_2$ \& helium gas at room temperature&(R)&$Z=3/8$\\\hline
(IV)&At very low pressure \& high temperature&(S)&$Z=1$\\\hline
&&(T)&Positive derivation from ideality\\\hline
&&(U)&Negative derivation from ideality\\\hline
\end{tabular}\end{center}
Select the CORRECT option -
\begin{qoptions}
\opt{A}{III $\to$ Q}
\opt{B}{II $\to$ S}
\opt{C}{IV $\to$ P}
\opt{D}{III $\to$ P}
\end{qoptions}
\qn{10} \begin{center}\footnotesize\setlength{\tabcolsep}{2pt}
\begin{tabular}{|L{0.8cm}L{2.2cm}|L{0.6cm}L{2.2cm}|}\hline
\multicolumn{2}{|L{3.0cm}|}{\centering List-I\\(condition of vander walls gas)}&\multicolumn{2}{L{2.8cm}|}{\centering List-II\\(compressibility factor)}\\\hline
(I)&At low to moderate pressure&(P)&$Z=1+\dfrac{Pb}{RT}$\\\hline
(II)&At critical condition of $P,V,T$&(Q)&$Z=1-\dfrac{a}{V_mRT}$\\\hline
(III)&For H$_2$ \& helium gas at room temperature&(R)&$Z=3/8$\\\hline
(IV)&At very low pressure \& high temperature&(S)&$Z=1$\\\hline
&&(T)&Positive derivation from ideality\\\hline
&&(U)&Negative derivation from ideality\\\hline
\end{tabular}\end{center}
Select the CORRECT option -
\begin{qoptions}
\opt{A}{I $\to$ P}
\opt{B}{II $\to$ R}
\opt{C}{IV $\to$ T}
\opt{D}{IV $\to$ U}
\end{qoptions}
\qn{11} \begin{center}\footnotesize\setlength{\tabcolsep}{2pt}
\begin{tabular}{|L{0.8cm}L{1.3cm}|L{0.6cm}L{3.9cm}|}\hline
\multicolumn{2}{|L{2.1cm}|}{\centering List-I\\(Molecules)}&\multicolumn{2}{L{4.5cm}|}{\centering List-II\\(Various facts)}\\\hline
(I)&Neo--C$_5$H$_{12}$&(P)&Maximum number of same atoms in one plane $=3$\\\hline
(II)&C$_2$H$_2$Cl$_2$&(Q)&$n$-isomer has more surface area hence more BP then iso -- and neo --\\\hline
(III)&CH$_3$Cl&(R)&Cis-form has more dipole moment than trans-form\\\hline
(IV)&CH$_4$&(S)&has more boiling point than tetrafluoromethane due to more repulsive forces in latter\\\hline
&&(T)&has more dipole moment than dichloromethane and trichloromethane\\\hline
&&(U)&have regular geometry around carbon atom(s)\\\hline
\end{tabular}\end{center}
Select the incorrect match.
\begin{qoptions}
\opt{A}{I $\to$ P}
\opt{B}{II $\to$ R}
\opt{C}{III $\to$ T}
\opt{D}{IV $\to$ S}
\end{qoptions}
\qn{12} \begin{center}\footnotesize\setlength{\tabcolsep}{2pt}
\begin{tabular}{|L{0.8cm}L{1.3cm}|L{0.6cm}L{3.9cm}|}\hline
\multicolumn{2}{|L{2.1cm}|}{\centering List-I\\(Molecules)}&\multicolumn{2}{L{4.5cm}|}{\centering List-II\\(Various facts)}\\\hline
(I)&Neo--C$_5$H$_{12}$&(P)&Maximum number of same atoms in one plane $=3$\\\hline
(II)&C$_2$H$_2$Cl$_2$&(Q)&$n$-isomer has more surface area hence more BP then iso -- and neo --\\\hline
(III)&CH$_3$Cl&(R)&Cis-form has more dipole moment than trans-form\\\hline
(IV)&CH$_4$&(S)&has more boiling point than tetrafluoromethane due to more repulsive forces in latter\\\hline
&&(T)&has more dipole moment than dichloromethane and trichloromethane\\\hline
&&(U)&have regular geometry around carbon atom(s)\\\hline
\end{tabular}\end{center}
Select the correct match.
\begin{qoptions}
\opt{A}{I $\to$ U}
\opt{B}{II $\to$ P}
\opt{C}{III $\to$ P}
\opt{D}{IV $\to$ T}
\end{qoptions}
\sect{SECTION-II}
\qn{1} 10 g impure mixture of KIO$_3$ was titrated with excess of KI in acidic medium. Iodine so liberated required 100 ml of 0.15 N solution of sodium thiosulphate for titration. The mass-percentage of KIO$_3$ in the mixture is \underline{\hspace{1.5cm}} .

(Atomic mass : K = 39, I = 127, O = 16)
\qn{2} The edge length of unit cell of a metal having molecular mass 75 g/mole is 4 \AA. If metal crystallises in cubic lattice and its density is 3.91 g/cc then find radius (in pm) of metal atom.

Given : $N_A=6\times10^{23}$, $\sqrt2=1.414$, $\sqrt3=1.732$
\qn{3} Consider the following molecules

CO$_2$ , BF$_3$ , H$_2$O , NH$_3$ , CH$_4$ , PCl$_5$ , SF$_6$ , IF$_7$ , PCl$_3$F$_2$ , PCl$_2$F$_3$ , XeF$_5^\ominus$ , ClF$_3$ , SF$_4$

$\dfrac{\text{number of molecules having all equal bond length}}{\text{number of polar molecules showing H - bonding}}=x$

Find $x$?
\qn{4} Regarding lewis dot structure of sulphate ion, consider the following statements

(a) Resonance is possible with bond order $=1.5$

(b) Each atom contains formal charge

(c) The expanded structure and lewis structure have equal number of S--O linkage

(d) Magnitude of formal charge on sulphur is equal to $\pi$ bonds in most stable expanded structure

(e) Ratio of total number of lone pairs to bond pairs is 13 : 4

(f) Formal charge of sulphur here is different from formal charge of sulphur in lewis dot structure of SO$_3$ and SO$_2$Cl$_2$

$\dfrac{\text{number of correct statements}}{\text{number of wrong statements}}=x$

Find $x$
\qn{5} In carius method of estimation of halogen 0.15 g of an organic compound gave 0.12 g of AgBr. Find out the percentage of bromine in the compound?
\qn{6} On complete combustion, 0.246 g of an organic compound gave 0.198 g of carbon dioxide and 0.1014g of water. Determine the percentage composition of carbon in the compound
\partsec{3}{MATHEMATICS}
\sect{SECTION-I(i)}
\qn{1} Consider the curve

$S=(\sin\theta)x^2+2(\cos\theta)xy+(\sin\theta)y^2+2(\cos\theta)x+2(\cos\theta)y+\sin\theta=0$

Now, choose the correct statement(s)
\begin{qoptions}
\opt{A}{If $\theta=\dfrac{3\pi}4$, then $S=0$ represents a pair of straight lines.}
\opt{B}{If $S=0$ represents a pair of real straight lines, then $\theta=n\pi+\dfrac\pi4$, $n\in Z$}
\opt{C}{If $S=0$ represents a pair of real straight lines, then they are parallel.}
\opt{D}{There exist some $\theta$ for which $S=0$ represents no real locus.}
\end{qoptions}
\qn{2} Let $A,B,C$ be 3 sets such that

$A=\left\{(x,y):\dfrac x{\cos\theta}=\dfrac y{\sin\theta}=2,\theta\in R\right\}$

$B=\left\{(x,y):\dfrac{x+1}{\cos2205^\circ}=\dfrac y{\sin2205^\circ}=r,r\in R\right\}$

$C=\left\{(x,y):(x-\sqrt3)^2+(y-1)^2\leqslant a^2;a\in R\right\}$

Let $C_1,C_2$ be the curves represented by $A$ and $B$ respectively while $R$ be the region represented by $C$ then which of following statement is/are correct
\begin{qoptions}
\opt{A}{A points $P(x_1,y_1)$ and $Q(x_2,y_2)$ be such that $P,Q\in A\cap B$, then point of intersection of tangents drawn to $C_1$ at $P$ and $Q$ is $(2,3)$.}
\opt{B}{A points $P(x_1,y_1)$ and $Q(x_2,y_2)$ be such that $P,Q\in A\cap B$, then point of intersection of tangents drawn to $C_1$ at $P$ and $Q$ is $(-4,4)$.}
\opt{C}{Minimum value of `$a$' for which $A\cap C=A$ is $\alpha$, then length of curve $C_2$ lying in $R$ corresponding to $a=\alpha$, is $\sqrt{58}$}
\opt{D}{Minimum value of `$a$' for which $A\cap C=A$ is $\alpha$, then length of curve $C_2$ lying in $R$ corresponding to $a=\alpha$, is $\sqrt{57}$}
\end{qoptions}
\qn{3} The sum of $99^{\text{th}}$ and $100^{\text{th}}$ term of the sequence 1, 2, 4, 6, 7, 10, 10, 14, ..... is divisible by
\begin{qoptions}
\opt{A}{2}
\opt{B}{53}
\opt{C}{109}
\opt{D}{173}
\end{qoptions}
\qn{4} In a game a person $P$ at point $(x,y)$ can jump to any one of the following points relatively: $(x+1,y)$, $(x+2,y)$, $(x,y+1)$ or $(x,y+2)$. If $P$ is at $(0,0)$, then he can reach
\begin{qoptions}
\opt{A}{$(2,2)$ in 14 ways}
\opt{B}{$(2,3)$ in 32 ways}
\opt{C}{$(2,2)$ in 9 ways}
\opt{D}{$(2,3)$ in 26 ways}
\end{qoptions}
\qn{5} Let $x=(7+4\sqrt3)^n$, where $n$ is a natural number, then -
\begin{qoptions}
\opt{A}{$[x]$ is an even integer (where $[.]$ denotes greatest integer function)}
\opt{B}{$[x]$ is an odd integer (where $[.]$ denotes greatest integer function)}
\opt{C}{$(1-\{x\})(x)=1$ (where $\{.\}$ denotes fractional part function)}
\opt{D}{$x\{x\}=1$ (where $\{.\}$ denotes fractional part function)}
\end{qoptions}
\qn{6} The set of all values of $x$ in $[0,\pi]$ for which the function $f(x)=\sqrt{1-2\sin3x}$ takes real values is-
\begin{qoptions}
\opt{A}{$\left[0,\dfrac\pi6\right]\cup\left[\dfrac{5\pi}6,\pi\right]$}
\opt{B}{$\left[0,\dfrac\pi9\right]\cup\left[\dfrac{2\pi}9,\dfrac{7\pi}9\right]\cup\left[\dfrac{8\pi}9,\pi\right]$}
\opt{C}{$\left[0,\dfrac\pi{18}\right]\cup\left[\dfrac{5\pi}{18},\dfrac{13\pi}{18}\right]\cup\left[\dfrac{17\pi}{18},\pi\right]$}
\opt{D}{$\left[0,\dfrac{13\pi}{18}\right]\cup\left[\dfrac{17\pi}{18},\pi\right]$}
\end{qoptions}
\qn{7} Sum of the solutions of the equation $[x^2]-2x+1=0$ is (where $[.]$ denotes greatest integer function)
\begin{qoptions}
\opt{A}{$\dfrac12$}
\opt{B}{2}
\opt{C}{3}
\opt{D}{$\dfrac32$}
\end{qoptions}
\qn{8} The range of the function $f(x)=\dfrac{\sqrt{1-x^2}}{1+|x|}$ is-
\begin{qoptions}
\opt{A}{$[0,1]$}
\opt{B}{$\left[0,\dfrac1{\sqrt2}\right]$}
\opt{C}{$\left[0,\dfrac12\right]$}
\opt{D}{$\left[0,\dfrac{\sqrt3}2\right]$}
\end{qoptions}
\sect{SECTION-I(ii)}
\qn{9} Answer the following by appropriately matching the lists based on the information given in the paragraph

$\alpha_1,\alpha_2,\alpha_3,\alpha_4$ and $\alpha_5$ are all roots of $x^5-4x^3-3x^2+2x-2=0$.

There are some expression given in List-I whose values are given in List-II below :
\begin{center}\footnotesize\setlength{\tabcolsep}{2pt}
\begin{tabular}{|L{0.8cm}L{4.4cm}|L{0.6cm}L{0.6cm}|}\hline
\multicolumn{2}{|c|}{List-I}&\multicolumn{2}{c|}{List-II}\\\hline
(I)&$\alpha_1^2+\alpha_2^2+\alpha_3^2+\alpha_4^2+\alpha_5^2$&(P)&1\\\hline
(II)&$\left|(1-\alpha_1^2)(1-\alpha_2^2)\dots(1-\alpha_5^2)\right|$&(Q)&8\\\hline
(III)&$\bigl|(\alpha_1^3-4\alpha_1-3)(\alpha_2^3-4\alpha_2-3)$ $\dots(\alpha_5^3-4\alpha_5-3)\bigr|$&(R)&12\\\hline
(IV)&$\dfrac1{\alpha_1}+\dfrac1{\alpha_2}+\dots\dfrac1{\alpha_5}$&(S)&24\\\hline
&&(T)&36\\\hline
&&(U)&48\\\hline
\end{tabular}\end{center}
Which of the following is the only correct combination ?
\begin{qoptions}
\opt{A}{I $\to$ R}
\opt{B}{II $\to$ S}
\opt{C}{III $\to$ T}
\opt{D}{IV $\to$ Q}
\end{qoptions}
\qn{10} Answer the following by appropriately matching the lists based on the information given in the paragraph

$\alpha_1,\alpha_2,\alpha_3,\alpha_4$ and $\alpha_5$ are all roots of $x^5-4x^3-3x^2+2x-2=0$.

There are some expression given in List-I whose values are given in List-II below :
\begin{center}\footnotesize\setlength{\tabcolsep}{2pt}
\begin{tabular}{|L{0.8cm}L{4.4cm}|L{0.6cm}L{0.6cm}|}\hline
\multicolumn{2}{|c|}{List-I}&\multicolumn{2}{c|}{List-II}\\\hline
(I)&$\alpha_1^2+\alpha_2^2+\alpha_3^2+\alpha_4^2+\alpha_5^2$&(P)&1\\\hline
(II)&$\left|(1-\alpha_1^2)(1-\alpha_2^2)\dots(1-\alpha_5^2)\right|$&(Q)&8\\\hline
(III)&$\bigl|(\alpha_1^3-4\alpha_1-3)(\alpha_2^3-4\alpha_2-3)$ $\dots(\alpha_5^3-4\alpha_5-3)\bigr|$&(R)&12\\\hline
(IV)&$\dfrac1{\alpha_1}+\dfrac1{\alpha_2}+\dots\dfrac1{\alpha_5}$&(S)&24\\\hline
&&(T)&36\\\hline
&&(U)&48\\\hline
\end{tabular}\end{center}
Which of the following is the only correct combination ?
\begin{qoptions}
\opt{A}{I $\to$ R}
\opt{B}{II $\to$ T}
\opt{C}{III $\to$ U}
\opt{D}{IV $\to$ R}
\end{qoptions}
\qn{11} Answer the following by appropriately matching the lists based on the information given in the paragraph

List-I contains some trigonometric equations and List-II provides number of solutions in given interval.
\begin{center}\footnotesize\setlength{\tabcolsep}{2pt}
\begin{tabular}{|L{0.8cm}L{3.8cm}|L{0.6cm}L{1.9cm}|}\hline
\multicolumn{2}{|c|}{List-I}&\multicolumn{2}{c|}{List-II}\\\hline
(I)&$2\cot2x-3\cot3x=\tan2x$&(P)&0 in $\left(0,\dfrac\pi2\right)$\\\hline
(II)&$\sin^{10}x+\cos^{10}x=\dfrac{29}{16}\cos^42x$&(Q)&2 in $(0,\pi)$\\\hline
(III)&$\cos^42x+2\sin^22x=17(\cos x+\sin x)^8$&(R)&4 in $(0,\pi)$\\\hline
(IV)&$\sin\!\left(x+\dfrac\pi9\right)=\cos\!\left(2x-\dfrac{59\pi}{180}\right)$&(S)&4 in $(0,2\pi)$\\\hline
&&(T)&2 in $\left(-\dfrac\pi2,\dfrac\pi2\right)$\\\hline
&&(U)&2 in $\left(\dfrac\pi2,\pi\right)$\\\hline
\end{tabular}\end{center}
Which of the following is the only correct combination ?
\begin{qoptions}
\opt{A}{I $\to$ T}
\opt{B}{II $\to$ R}
\opt{C}{III $\to$ P}
\opt{D}{IV $\to$ R}
\end{qoptions}
\qn{12} Answer the following by appropriately matching the lists based on the information given in the paragraph

List-I contains some trigonometric equations and List-II provides number of solutions in given interval.
\begin{center}\footnotesize\setlength{\tabcolsep}{2pt}
\begin{tabular}{|L{0.8cm}L{3.8cm}|L{0.6cm}L{1.9cm}|}\hline
\multicolumn{2}{|c|}{List-I}&\multicolumn{2}{c|}{List-II}\\\hline
(I)&$2\cot2x-3\cot3x=\tan2x$&(P)&0 in $\left(0,\dfrac\pi2\right)$\\\hline
(II)&$\sin^{10}x+\cos^{10}x=\dfrac{29}{16}\cos^42x$&(Q)&2 in $(0,\pi)$\\\hline
(III)&$\cos^42x+2\sin^22x=17(\cos x+\sin x)^8$&(R)&4 in $(0,\pi)$\\\hline
(IV)&$\sin\!\left(x+\dfrac\pi9\right)=\cos\!\left(2x-\dfrac{59\pi}{180}\right)$&(S)&4 in $(0,2\pi)$\\\hline
&&(T)&2 in $\left(-\dfrac\pi2,\dfrac\pi2\right)$\\\hline
&&(U)&2 in $\left(\dfrac\pi2,\pi\right)$\\\hline
\end{tabular}\end{center}
Which of the following is the only correct combination ?
\begin{qoptions}
\opt{A}{I $\to$ R}
\opt{B}{II $\to$ T}
\opt{C}{III $\to$ U}
\opt{D}{IV $\to$ R}
\end{qoptions}
\sect{SECTION-II}
\qn{1} If a normal chord $AB$ to the hyperbola $xy=c^2$, normal at $A$, subtends an angle $\theta$ at centre and $\angle OAB=\alpha$, ($O$ is origin) then $\left|\dfrac{\tan\alpha}{\tan\theta}\right|$ is
\qn{2} Let $P$ be any point on ellipse $3x^2+4y^2=12$ and $S,S'$ are its foci then the locus of the centroid of triangle $PSS'$ is a conic $C$. Area of quadrilateral formed by tangents at the extremities of latus-rectum of the conic $C$ equals
\qn{3} If $y^2=4x$ and $y^2=(x-k)$ have a common normal other than $x$-axis, then the range of values of $k$ is $(p,\infty)$ ; where $p$ is equal to
\qn{4} From a point $P(0,5)$ on the circle $x^2+y^2=25$ two equal chords $PA$ and $PB$ of length 8 units are drawn. The equation of $AB$ is $y=k$, where $k$ is
\qn{5} The number of real solutions of the equation $\left(\dfrac{\log_{10}x}3\right)^{\log_{10}^2x-3\log_{10}x+3}=\log_{10}\sqrt[3]x$, is
\qn{6} The number of solutions of the equation $\cos^{12}x+3\cos^4x\sin^3x=\sin^9x-1$ in $x\in[0,100\pi]$ is
\end{multicols}
\end{document}
